\chapter{Zatrzymywanie, nierówności maksymalne i zbieżność martyngałów}
\label{ch:stop-cont}

W całym rozdziale zakładamy warunki zwykłe filtracji,
jak w rozdziale~\ref{ch:mart-cont}. Dzięki zupełności
możemy wybierać wszędzie c\`adl\`ag trajektorie
oraz poprawiać czasy zatrzymania na zbiorach zerowych.

\section{Co wolno zrobić w losowej chwili?}
Jeżeli znamy przyszłość, możemy zawsze „wybrać” korzystną
chwilę sprzedaży. Twierdzenia o zatrzymywaniu wyjaśniają,
dlaczego nie uzyskujemy takiego przywileju, gdy chwilę
zatrzymania wybieramy bez zaglądania w przyszłość.
Samo pojęcie czasu zatrzymania nie wystarcza jednak do
przeniesienia równości wartości oczekiwanych na nieograniczony
horyzont. W każdym zastosowaniu trzeba kontrolować
całkowalność rodziny zatrzymanych wartości.

\begin{twierdzenie}[Próbkowanie opcjonalne w ograniczonym horyzoncie]
\label{thm:ost-cont}
Niech $(\calF_t)$ będzie filtracją spełniającą
warunki zwykłe, $M$ martyngałem
c\`adl\`ag i $\sigma,\tau$ czasami zatrzymania takimi,
że $0\le\sigma\le\tau\le T$ prawie na pewno dla ustalonego
$T<\infty$. Wtedy $M_\sigma,M_\tau\in L^1$ i
\[
\E[M_\tau\mid\calF_\sigma]=M_\sigma
\qquad\text{prawie na pewno}.
\]
Dla podmartyngału (odpowiednio nadmartyngału) c\`adl\`ag
ten sam wniosek zachodzi z odpowiednią nierównością,
jeżeli $\{M_\rho:\rho\le T\text{ jest czasem zatrzymania}\}$
jest rodziną jednostajnie całkowalną (w szczególności
gdy proces jest klasy $(DL)$).
\end{twierdzenie}
\begin{proof}
Przypadek $T=0$ jest natychmiastowy.
Gdy ograniczenia i uporządkowanie czasów zachodzą
tylko p.n., na wspólnym zbiorze zerowym przyjmujemy
oba czasy równe zero. Zbiór ten należy do $\calF_0$,
więc zmienione czasy nadal są czasami zatrzymania.
Ich $\sigma$-ciała są identyczne z pierwotnymi:
zdarzenia testujące definicję $\calF_\sigma$
zmieniają się tylko o podzbiory tego zbioru zerowego,
należące do każdej $\calF_t$.
Wartości w czasach zatrzymania są mierzalne
w odpowiednich $\sigma$-ciałach na mocy
twierdzenia~\ref{thm:sampled}.
Możemy zatem zakładać nierówności punktowe.
Podzielmy odcinek $[0,T]$ na $2^n$ równych części
o końcach $t_k^n=kT/2^n$. Niech $\sigma_n,\tau_n$
będą najmniejszymi punktami tej siatki, nie mniejszymi
od odpowiednich czasów (dla wartości $T$ przyjmujemy $T$).
Wówczas $\sigma_n,\tau_n$ są czasami zatrzymania,
$\sigma_n\le\tau_n\le T$ i
$\sigma_n\downarrow\sigma$, $\tau_n\downarrow\tau$.
Istotnie zdarzenie $\{\sigma_n\le t\}$ jest
$\{\sigma\le t_k^n\}$, gdzie $t_k^n$ jest ostatnim
punktem siatki nie większym od $t$ (dla $t\ge T$
zdarzenie jest całym $\Omega$); analogicznie dla $\tau_n$.
Zagnieżdżenie siatek daje monotoniczność przybliżeń.
W skończonym schemacie próbkowania $\calF_{\sigma_n}$
jest dokładnie $\sigma$-ciałem dyskretnego czasu:
dla każdej wartości $t_k^n$ zdarzenie
$A\cap\{\sigma_n=t_k^n\}$ musi należeć
do $\calF_{t_k^n}$. Równoważność tej charakteryzacji
z definicją wynika przez rozkład na skończenie wiele
wartości czasu.
Wartości $M_{t_k^n}$ tworzą martyngał dyskretny, więc
jego twierdzenie o zatrzymaniu daje dla
$A\in\calF_\sigma\subseteq\calF_{\sigma_n}$
\[
\E[M_{\tau_n}\chi_A]=\E[M_{\sigma_n}\chi_A].
\]
Ciągłość prawostronna trajektorii zapewnia zbieżność
$M_{\tau_n}\to M_\tau$ i $M_{\sigma_n}\to M_\sigma$
prawie na pewno. Żeby uzasadnić przejście do wartości
oczekiwanych, zauważmy, że w skończonym schemacie próbkowania
$M_{\tau_n}=\E[M_T\mid\calF_{\tau_n}]$ oraz
$M_{\sigma_n}=\E[M_T\mid\calF_{\sigma_n}]$.
Rodziny te są jednostajnie całkowalne na mocy
lematu~\ref{lem:mart-cont-ui-cond}. Ich zbieżność prawie na pewno
jest zatem także zbieżnością w $L^1$.
Przejście $n\to\infty$ daje
$\E[M_\tau\chi_A]=\E[M_\sigma\chi_A]$ dla każdego
$A\in\calF_\sigma$, czyli żądaną równość warunkową.
W przypadku podmartyngałów i nadmartyngałów
zakładana jednostajna całkowalność wartości w ograniczonych
czasach zatrzymania pozwala przejść do granicy w $L^1$
dokładnie w ten sam sposób, przy zachowaniu właściwego
zwrotu nierówności. Dla ogólnych procesów bez tej
kontroli nie uzasadniamy przejścia do granicy przez
samą zbieżność prawie na pewno.
\end{proof}

\begin{uwaga}[Związek z procesem zatrzymanym]
Każdy martyngał c\`adl\`ag jest klasy $(DL)$
(twierdzenie~\ref{thm:mart-cont-class-d}).
Dla podmartyngału założenie o jednostajnej
całkowalności wartości w czasach zatrzymania
jest warunkiem dostatecznym powyższego dowodu,
a nie definicją podmartyngału.
Twierdzenie o zatrzymywaniu procesu
(twierdzenie~\ref{thm:mart-cont-stopping})
jest zastosowaniem ograniczonego próbkowania
opcjonalnego; nie powtarzamy jego dowodu.
\end{uwaga}

\begin{zadanie}[Nie wolno podejmować decyzji na podstawie przyszłości]
Niech $W$ będzie standardowym procesem Wienera
w uzupełnionej filtracji naturalnej i określmy
losową chwilę $\rho=0$ na $\{W_1\ge0\}$ oraz $\rho=1$
na $\{W_1<0\}$. Czy $\rho$ jest czasem zatrzymania?
Ile wynosi $\E W_\rho$?

\emph{Rozwiązanie.} Zdarzenie $\{\rho\le0\}=\{W_1\ge0\}$
nie należy do $\calF_0$ uzupełnionej filtracji naturalnej
ruchu Browna: ma prawdopodobieństwo $1/2$ i zależy od przyszłej
obserwacji. Zatem $\rho$ nie jest czasem zatrzymania.
Ponieważ $W_0=0$ i $W_1\sim\Normal(0,1)$,
\[
\E W_\rho=\E[W_1\chi_{\{W_1<0\}}]
=-\frac1{\sqrt{2\pi}},
\qquad
\text{gdyż }\int_{-\infty}^0 x e^{-x^2/2}\,
\mathrm d\lambda_1(x)=-1.
\]
Nierówność średnich nie przeczy twierdzeniu o próbkowaniu
opcjonalnym: jego podstawowe założenie, czyli brak
korzystania z przyszłości, nie jest spełnione.
\end{zadanie}


\section{Maksimum procesu a jego wartość końcowa}
\begin{twierdzenie}[Nierówność maksymalna Dooba]
\label{thm:doob-cont}
Niech $X$ będzie nieujemnym podmartyngałem c\`adl\`ag
na $[0,T]$ i $X_T^*=\sup_{0\le t\le T}X_t$.
Dla $a>0$ zachodzi nierówność ogonowa
\[
a\Prob(X_T^*\ge a)\le
\E[X_T\chi_{\{X_T^*\ge a\}}]\le\E X_T.
\]
Jeżeli $p>1$ i $X_T\in L^p$, to
\[
\E[(X_T^*)^p]\le
\left(\frac{p}{p-1}\right)^p\E X_T^p.
\]
W szczególności dla martyngału c\`adl\`ag $M$
można zastosować obie nierówności do $X_t=|M_t|$.
\end{twierdzenie}
\begin{proof}
Ustalmy skończony zbiór chwil
$0=t_0<\cdots<t_n=T$ i połóżmy $Y_k=X_{t_k}$.
Jest to nieujemny podmartyngał dyskretny.
Niech $Y=\max_{0\le k\le n}Y_k$.
Jeżeli $\rho=\min\{k:Y_k\ge a\}\wedge n$,
z konwencją $\inf\varnothing=\infty$,
to na zdarzeniu $\{Y\ge a\}$ mamy $Y_\rho\ge a$.
Próbkowanie podmartyngału daje
$a\Prob(Y\ge a)\le\E[Y_n\chi_{\{Y\ge a\}}]$:
dokładniej sumujemy nierówności warunkowe po rozłącznych
zdarzeniach $\{\rho=k,Y\ge a\}$, które należą
do $\calF_{t_k}$: dla $k<n$ jest to $\{\rho=k\}$,
a dla $k=n$ przecinamy $\{\rho=n\}$ z $\{Y_n\ge a\}$.
Dla $p>1$ z $0\le Y_k\le\E[X_T\mid\calF_{t_k}]$
p.n. i warunkowej nierówności Jensena wynika
$\E Y_k^p\le\E X_T^p$. W szczególności
$\E Y^p\le\sum_{k=0}^n\E Y_k^p<\infty$.
Całkujemy nierówność ogonową względem
$p a^{p-2}\,\mathrm d\lambda_1(a)$ po $a\in(0,\infty)$.
Twierdzenie Tonellego daje
\[
\E Y^p\le\frac{p}{p-1}\E[Y_nY^{p-1}].
\]
Z nierówności H\"oldera prawa strona jest nie większa od
$\frac{p}{p-1}\|Y_n\|_p\|Y\|_p^{p-1}$.
Jeśli $\E Y^p=0$, teza jest oczywista;
w przeciwnym razie dzielimy przez odpowiednią potęgę
i otrzymujemy $\|Y\|_p\le p(p-1)^{-1}\|Y_n\|_p$.
Przechodzimy do rosnącego ciągu skończonych siatek,
których suma jest gęsta i zawiera $T$.
Dzięki prawostronnej ciągłości trajektorii
$Y$ rośnie do $X_T^*$, który jest zatem mierzalny.
Zbieżność monotoniczna daje oszacowanie w $L^p$.
W wersji ogonowej supremum funkcji c\`adl\`ag
może nie być osiągane. Dla $0<\varepsilon<a$
zdarzenia $\{Y\ge a-\varepsilon\}$ rosną wraz
z zagęszczaniem siatek do zdarzenia $A_\varepsilon$,
które zawiera $\{X_T^*\ge a\}$ i jest zawarte
w $\{X_T^*\ge a-\varepsilon\}$.
Dyskretna nierówność i zbieżność monotoniczna
dają $(a-\varepsilon)\Prob(A_\varepsilon)
\le\E[X_T\chi_{A_\varepsilon}]$.
Dla ciągu $\varepsilon\downarrow0$ zdarzenia
$A_\varepsilon$ maleją do $\{X_T^*\ge a\}$,
więc przejście do granicy daje nierówność ogonową.
Dla $X=|M|$ używamy warunkowej nierówności Jensena,
która dowodzi własności podmartyngałowej $X$.
\end{proof}

\begin{wniosek}[Maksimum ruchu Browna w $L^2$]
Dla standardowego procesu Wienera $W$ oraz $T>0$
\[
\E\sup_{0\le t\le T}|W_t|^2\le4T.
\]
Dla standardowego $d$-wymiarowego ruchu Browna $B$
\[
\E\sup_{0\le t\le T}|B_t|^2\le4dT.
\]
\end{wniosek}
\begin{proof}
Pierwsza nierówność wynika z twierdzenia dla $p=2$
i $\E W_T^2=T$. W drugiej używamy
$|B_t|^2=\sum_i(B_t^{(i)})^2$ oraz
$\sup_t\sum_i a_i(t)\le\sum_i\sup_t a_i(t)$.
\end{proof}

\section{Nierówność przekroczeń i istnienie granic}
\begin{definicja}[Przekroczenie przedziału]
Dla funkcji rzeczywistej $x$ i liczb $a<b$ \emph{
pełne przekroczenie w górę przedziału $[a,b]$} to
para chwil $u<v$ takich, że $x(u)\le a$
oraz $x(v)\ge b$, przy czym po osiągnięciu $b$
rozpoczynamy poszukiwanie następnego przekroczenia
od ponownego zejścia do poziomu $a$. Maksymalną liczbę
takich rozłącznych przekroczeń do czasu $T$ oznaczamy
$U_{a,b}(T)$ (supremum liczby par w skończonych
ciągach $u_1<v_1<u_2<v_2<\cdots\le T$).
W dowodach używamy liczby przekroczeń na ustalonych
przeliczalnych zbiorach chwil, mierzalnej jako
supremum po skończonych podzbiorach.
\end{definicja}

\begin{lemat}[Przekroczenia dla martyngału dyskretnego]
Niech $(M_k)_{k=0}^n$ będzie martyngałem całkowalnym,
a $a<b$. Wtedy
\[
(b-a)\E U_{a,b}(n)\le\E[(a-M_n)^+].
\]
\end{lemat}
\begin{proof}
Budujemy przewidywalną strategię: nabywamy jedną jednostkę
aktywa, gdy $M_k\le a$, sprzedajemy przy pierwszym
$M_l\ge b$ i ponawiamy ten schemat. Pozycję otwartą
w chwili $n$ wyceniamy według $M_n$.
Dla każdego pełnego cyklu zysk jest co najmniej $b-a$.
Ewentualna niedokończona pozycja otwarta po cenie
co najwyżej $a$ ma zysk większy lub równy
$M_n-a\ge-(a-M_n)^+$. Łączny zysk $G_n$ spełnia
$G_n\ge(b-a)U_{a,b}(n)-(a-M_n)^+$.
Aby uzasadnić przewidywalność, definiujemy $H_k$ w chwili $k$
wyłącznie na podstawie $M_0,\ldots,M_k$: po zakupie utrzymujemy
$H_k=1$ do pierwszego osiągnięcia poziomu $b$, a po sprzedaży
przyjmujemy $H_k=0$ aż do kolejnego osiągnięcia poziomu $a$.
W szczególności $H_k\in\{0,1\}$ jest $\calF_k$-mierzalne.
Ponieważ
$G_n=\sum_{k=0}^{n-1}H_k(M_{k+1}-M_k)$ jest sumą skończoną,
$\E|G_n|<\infty$ i
$\E G_n=\sum_k\E[H_k\E(M_{k+1}-M_k\mid\calF_k)]=0$. Biorąc wartości
oczekiwane w nierówności, otrzymujemy tezę.
\end{proof}

\begin{twierdzenie}[Zbieżność przy ograniczeniu w $L^1$]
\label{thm:mart-cont-convergence}
Niech $M$ będzie podmartyngałem c\`adl\`ag na
$[0,\infty)$ takim, że $\sup_{t\ge0}\E M_t^+<\infty$.
Wtedy $\sup_t\E|M_t|<\infty$ oraz istnieje zmienna
losowa $M_\infty\in L^1$ taka, że
$M_t\to M_\infty$ prawie na pewno przy $t\to\infty$.
W szczególności wynik dotyczy martyngałów ograniczonych
w $L^1$. Dla nadmartyngału wystarcza
$\sup_t\E M_t^-<\infty$.
Nie wnioskujemy z tych założeń zbieżności w $L^1$.
\end{twierdzenie}
\begin{proof}
Z monotoniczności średnich podmartyngału mamy
$\E|M_t|=2\E M_t^+-\E M_t\le2\sup_u\E M_u^+-\E M_0$.
Wybierzmy przeliczalny gęsty zbiór chwil
$D=\Q\cap[0,\infty)$ oraz rosnące skończone
podzbiory zawierające zero i wyczerpujące $D$.
Na każdej siatce próbkowany proces jest
podmartyngałem dyskretnym. Nierówność przekroczeń
Dooba (tom I, rozdział~\ref{I-ch:nadmartyngaly})
daje dla ostatniej chwili $v$ siatki
$(b-a)\E U\le\E[(M_v-a)^+]-\E[(M_0-a)^+]
\le\sup_t\E M_t^++|a|$.
Zbieżność monotoniczna oraz przeliczalność
par wymiernych $a<b$ dają wspólny zbiór pełnej
miary, na którym liczba przekroczeń każdej
takiej pary na $D$ jest skończona.
Gdyby wartości $M_t$ dla $t\in D$ nie miały granicy
w rozszerzonej prostej rzeczywistej przy $t\to\infty$,
ich dolna granica byłaby ściśle mniejsza od górnej.
Wybierając pomiędzy nimi wymierne $a<b$,
moglibyśmy naprzemiennie wybierać coraz późniejsze
chwile z wartościami poniżej $a$ i powyżej $b$,
uzyskując nieskończenie wiele przekroczeń.
Zatem granica istnieje.
Stosując lemat Fatou wzdłuż całkowitych chwil
$n\to\infty$, dostajemy
$\E\liminf_n |M_n|\le\sup_t\E|M_t|<\infty$.
Ponieważ granica wzdłuż $D$ już istnieje p.n.,
jest skończona i należy do $L^1$.
Jako mierzalną wersję $M_\infty$ bierzemy granicę
$(M_n)$ tam, gdzie jest skończona, i zero poza tym;
jest ona $\calF_\infty$-mierzalna, gdzie
$\calF_\infty=\sigma(\bigcup_t\calF_t)$.
Prawostronna ciągłość trajektorii rozszerza wynik
na wszystkie chwile: dla każdej chwili $t$
można dobrać $r\in D$ tuż za $t$ z
$|M_r-M_t|$ dowolnie małym. Gdyby pewien ciąg
$t_n\to\infty$ nie dawał tej samej granicy,
dobierając $r_n\ge t_n$ z błędem nie większym niż
$1/n$, otrzymalibyśmy sprzeczność ze zbieżnością
próbek wymiernych.
Dla nadmartyngału stosujemy dowód do $-M$.
\end{proof}

\begin{twierdzenie}[Zbieżność jednostajnie całkowalnego martyngału]
\label{thm:mart-cont-ui-convergence}
Niech $M$ będzie martyngałem c\`adl\`ag na
$[0,\infty)$, a rodzina $\{M_t:t\ge0\}$
jednostajnie całkowalna. Wtedy istnieje
$M_\infty\in L^1$ takie, że $M_t\to M_\infty$
prawie na pewno i w $L^1$ oraz
\[
M_t=\E[M_\infty\mid\calF_t]
\qquad (t\ge0).
\]
Odwrotnie, każda taka reprezentacja z $M_\infty\in L^1$
tworzy jednostajnie całkowalny martyngał.
\end{twierdzenie}
\begin{proof}
Jednostajna całkowalność zapewnia ograniczenie w $L^1$,
więc istnieje granica prawie pewna z poprzedniego
twierdzenia. Kryterium Vitaliego daje też zbieżność
w $L^1$. Dla każdego ustalonego $s$ mamy
$M_s=\E[M_t\mid\calF_s]$ przy $t\ge s$.
Kontrakcja warunkowania w $L^1$ pozwala przejść
do granicy $t\to\infty$ po prawej stronie.
Wtedy $M_s=\E[M_\infty\mid\calF_s]$.
Implikacja odwrotna wynika z lematu~\ref{lem:mart-cont-ui-cond}
oraz własności wieży.
Związek z klasą $(D)$ podaje
twierdzenie~\ref{thm:mart-cont-class-d}.
\end{proof}

\begin{wniosek}[Zbieżność martyngałów w $L^p$]
\label{cor:mart-cont-lp-convergence}
Niech $p>1$, a $M$ będzie martyngałem c\`adl\`ag
spełniającym $\sup_{t\ge0}\E|M_t|^p<\infty$.
Wtedy istnieje $M_\infty\in L^p$ takie, że
$M_t\to M_\infty$ prawie na pewno i w $L^p$.
\end{wniosek}
\begin{proof}
Ograniczenie w $L^p$ implikuje jednostajną całkowalność
rodziny $(M_t)$, więc otrzymujemy granicę w $L^1$
i reprezentację $M_t=\E[M_\infty\mid\calF_t]$.
Lemat Fatou daje $M_\infty\in L^p$.
Warunkowa nierówność Jensena implikuje
$|M_t|^p\le\E[|M_\infty|^p\mid\calF_t]$.
Rodzina po prawej stronie jest jednostajnie całkowalna,
a więc taka jest również $\{|M_t|^p:t\ge0\}$.
Ponieważ $|M_t-M_\infty|^p\le2^{p-1}
(|M_t|^p+|M_\infty|^p)$, rodzina tych różnic do
potęgi $p$ też jest jednostajnie całkowalna.
Zbieżność prawie na pewno i kryterium Vitaliego dają
$\E|M_t-M_\infty|^p\to0$.
\end{proof}

\section{Kiedy można zatrzymać martyngał bez ograniczenia czasu?}
\begin{twierdzenie}[Próbkowanie dla jednostajnie całkowalnego martyngału]
\label{thm:mart-cont-ui-ost}
Niech $M$ będzie jednostajnie całkowalnym martyngałem
c\`adl\`ag, $\tau$ skończonym prawie na pewno czasem
zatrzymania. Wtedy $M_\tau\in L^1$ i
$\E M_\tau=\E M_0$. Ponadto, jeśli $\sigma\le\tau$
jest innym skończonym prawie na pewno czasem zatrzymania,
to $\E[M_\tau\mid\calF_\sigma]=M_\sigma$.
\end{twierdzenie}
\begin{proof}
Na mocy twierdzeń~\ref{thm:mart-cont-ui-convergence}
i~\ref{thm:mart-cont-class-d} dla każdego p.n.
skończonego czasu zatrzymania $\rho$ mamy
$M_\rho=\E[M_\infty\mid\calF_\rho]$.
Jeżeli $\sigma\le\tau$ p.n., po zmianie na zbiorze
zerowym możemy przyjąć uporządkowanie punktowe.
Wtedy $\calF_\sigma\subseteq\calF_\tau$,
a własność wieży daje
\[
\E[M_\tau\mid\calF_\sigma]
=\E[M_\infty\mid\calF_\sigma]=M_\sigma.
\]
Wybór $\sigma=0$ i wzięcie wartości oczekiwanej
dają $\E M_\tau=\E M_0$.
\end{proof}

\begin{zadanie}[Oczekiwanie na pierwszy sukces a uczciwa gra]
Niech $W$ będzie standardowym procesem Wienera,
$\tau_a=\inf\{t\ge0:W_t=a\}$, $a>0$. Wyjaśnić,
dlaczego nie możemy zastosować twierdzenia o
nieograniczonym zatrzymywaniu do $M=W$ i $\tau_a$.

\emph{Rozwiązanie.} Chociaż $\tau_a<\infty$ prawie
na pewno i $W_{\tau_a}=a$, martyngał $W$ nie jest
jednostajnie całkowalny na całej półprostej
($\E|W_t|$ rośnie jak $\sqrt t$).
Równość $\E W_{\tau_a}=\E W_0$ byłaby zatem
nieuprawniona. Ograniczone czasy $\tau_a\wedge T$
są dozwolone i dają $\E W_{\tau_a\wedge T}=0$;
brakuje jednak jednostajnej całkowalności przy
$T\to\infty$. Na tym właśnie polega różnica
między lokalną a globalną uczciwością.
\end{zadanie}

\begin{zadanie}[Ruina z dwiema granicami w czasie ciągłym]
Niech $W$ będzie standardowym procesem Wienera,
$a,b>0$ i
$\tau=\inf\{t\ge0:W_t\notin(-a,b)\}$.
Obliczyć prawdopodobieństwo wyjścia górą i wartość
oczekiwaną $\tau$, stosując tylko martyngały
$W_t$ oraz $W_t^2-t$.

\emph{Rozwiązanie.} Zatrzymujemy drugi martyngał w
$\tau\wedge n$:
$\E W_{\tau\wedge n}^2=\E(\tau\wedge n)$.
Ponieważ $|W_{\tau\wedge n}|\le\max(a,b)$,
$\E(\tau\wedge n)\le\max(a,b)^2$.
Zbieżność monotoniczna daje $\E\tau<\infty$, więc
$\tau<\infty$ prawie na pewno. Z ciągłości
$W_\tau\in\{-a,b\}$.
Proces $W^\tau$ jest ograniczony, a zatem
jednostajnie całkowalny. Z próbkowania opcjonalnego
$\E W_\tau=0$. Jeśli $p=\Prob(W_\tau=b)$, to
$bp-a(1-p)=0$, skąd
\[
\Prob(W_\tau=b)=\frac{a}{a+b},\qquad
\Prob(W_\tau=-a)=\frac{b}{a+b}.
\]
Ponownie używając $W_t^2-t$ i zbieżności zmajoryzowanej
po stronie kwadratu otrzymujemy
$\E\tau=\E W_\tau^2=b^2p+a^2(1-p)=ab$.
Rachunek ten pokazuje, jak zamienić zadanie o czasie
wyjścia na dwa równania dla momentów zatrzymanego
martyngału, bez rozwiązywania równania różniczkowego.
\end{zadanie}

\begin{zadanie}[Transformata Laplace'a czasu pierwszego trafienia]
Niech $a>0$ oraz
$\tau_a=\inf\{t\ge0:W_t=a\}$. Dla $u>0$
obliczyć $\E e^{-u\tau_a}$ bez posługiwania się
gęstością czasu trafienia.

\emph{Rozwiązanie.} Dla $\theta>0$ proces
$Z_t=\exp(\theta W_t-\theta^2t/2)$ jest
martyngałem. Dla każdego $T$ jego zatrzymana
wartość spełnia $0<Z_{T\wedge\tau_a}\le e^{\theta a}$:
przed trafieniem $W_s<a$, a składnik
$-\theta^2s/2$ jest niedodatni. Zatem rodzina
zatrzymanych wartości jest jednostajnie całkowalna.
Ponieważ $\tau_a<\infty$ prawie na pewno,
z próbkowania opcjonalnego i zbieżności zmajoryzowanej
\[
1=\E Z_{\tau_a}
=e^{\theta a}\E e^{-\theta^2\tau_a/2}.
\]
Po podstawieniu $\theta=\sqrt{2u}$ dostajemy
$\E e^{-u\tau_a}=e^{-a\sqrt{2u}}$.
Zauważmy, że tego zatrzymania \emph{nie} wolno
zastosować wprost do $W$, ponieważ $W_{\tau_a}$
ma średnią $a$, a $W_0$ średnią zero. Wykładnik
jest odpowiednim narzędziem właśnie dlatego, że
po zatrzymaniu posiada deterministyczne ograniczenie.
\end{zadanie}


% ================================================================
% ROZDZIAŁ 13. WARIACJA KWADRATOWA
% ================================================================
